\scanpage{20}\mspage{\unsure{17|19}}
% @ first three lines
to arrive at the notion of locus, and of interpretation of the data in terms of loci.

% @ terminology
\unsure{Point|Point} of view terminology, instead of ``elementary figure'', I now prefer to say ``\uline{multistratum}'', more telling.

Thus an arbitrary figure is a set of mutually compatible multistrata, closed under $\leq$, and locally finite.

% @ definition of locus
\deleted{\unsure{F|E}} We call a \deleted{multi} \uline{locus} of the country, a multistratum which is minimal for $\ll$, i.e. which has no refinement other than itself. \deleted{It is therefore} \margin{\deleted{\unsure{point|that} \unsure{of a point}}} \deleted{It is therefore} \deleted{In the} (In the typical examples, it is a set of subsets of $\deleted{\unsure{A|F}}X$, reduced to a single subset, which is punctual -- thus the loci are the points of $X$ \unsure{.-|..|.})

% @ minimal multistrata
These are therefore also multistrata minimal for $\leq$, \deleted{\unsure{therefore}} \inserted{a fortiori} ``manifolds'' \deleted{connected nonempty)''} -- these are the \unsure{0-connected|o-connected} non-manifolds which have no refinement other than itself.

% @ last lines and inclusion
Let $\mathcal{L}$ be the set of loci, \deleted{For every} thus
\[
\begin{aligned}
\mathcal{L}\subset M\quad (={}&\mathcal{F}_{el}=\text{set of}\\
&\hphantom{\mathcal{F}_{el}={}}\text{multistrata}\,)
\end{aligned}
\]
